\chapter{Metodología}\label{chap:metodologia}

% Barrera de flotantes: evita que las tablas migren de sección. Si el paquete
% placeins esta cargado se usa su \FloatBarrier; si no, se recurre a \clearpage.
\providecommand{\FloatBarrier}{\clearpage}

Este capítulo describe la metodología seguida para comparar de forma experimental los tres
\textit{frameworks} de aprendizaje federado seleccionados: Flower, NVIDIA FLARE y FedML. Se
presentan primero el entorno de trabajo y el equipo sobre el que se ejecutaron las pruebas, y a
continuación el marco de evaluación, que se articula en tres ejes: una comparativa cuantitativa
de rendimiento y dos comparativas cualitativas, de usabilidad y de flexibilidad. Para cada eje se
detallan las métricas, las rúbricas y los criterios empleados. El capítulo termina con las
medidas adoptadas para garantizar la reproducibilidad y el control de variables, y con el alcance
y las limitaciones del estudio. El escenario de entrenamiento común y la batería concreta de
experimentos se describen en el Capítulo~\ref{chap:experimentos}.

\section{Entorno de trabajo}\label{sec:entorno-trabajo}

Todos los experimentos se ejecutaron sobre el mismo equipo y bajo el mismo sistema operativo, de
modo que el entorno físico no introdujera diferencias entre las herramientas comparadas. Las
ejecuciones se realizaron sobre un equipo portátil con procesador Intel Core i7-13700H de 13.\textsuperscript{a}
generación (20 CPU lógicas), 14\,GiB de memoria RAM y una GPU NVIDIA GeForce RTX 3050 Laptop con
6\,GB de memoria de vídeo (VRAM), bajo Ubuntu 26.04 LTS. La GPU fue detectada correctamente por
PyTorch en los tres entornos.

Cada \textit{framework} se ejecutó, sin embargo, en un entorno Conda independiente, con versiones
distintas de Python, PyTorch y CUDA, debido a las dependencias propias de cada herramienta. La
Tabla~\ref{tab:entornos} recoge las versiones exactas empleadas. Esta diferencia se reconoce como
una amenaza menor a la validez de la comparación de rendimiento; se mitiga dejando registradas las
versiones de cada caso, lo que además favorece la reproducibilidad.

\begin{table}[htbp]
  \centering
  \small
  \caption{Versiones de software empleadas en cada entorno de ejecución.}
  \label{tab:entornos}
  \begin{tabular}{llllll}
    \toprule
    \textit{Framework} & Entorno & Python & Versión & PyTorch & CUDA \\
    \midrule
    Flower       & tfg-flower & 3.12.13 & 1.29.0 & 2.11.0+cu128 & 12.8 \\
    NVIDIA FLARE & nvflare    & 3.12.13 & 2.7.2  & 2.11.0+cu130 & 13.0 \\
    FedML        & fedml-tfg  & 3.10.20 & 0.9.6  & 2.12.0+cu126 & 12.6 \\
    \bottomrule
  \end{tabular}
\end{table}

\section{Comparativa de rendimiento: métricas e instrumentación}\label{sec:metricas-rendimiento}

La comparativa de rendimiento es la parte cuantitativa del estudio. Distingue entre métricas de
aprendizaje, que describen la calidad del modelo obtenido, y métricas de consumo de recursos, que
describen el coste de obtenerlo. Las métricas de aprendizaje son la precisión (\textit{accuracy})
y la pérdida (\textit{loss}) del modelo global sobre el conjunto de prueba centralizado, junto con
la pérdida media de entrenamiento. La evaluación sobre el conjunto de prueba se realiza de forma
externa al entrenamiento, cargando el modelo global final y midiendo su comportamiento sobre las
10\,000 imágenes reservadas, de modo idéntico para los tres \textit{frameworks}.

Las métricas de consumo de recursos se recogen durante toda la ejecución y se resumen en la
Tabla~\ref{tab:metricas}. Para la GPU se muestrea cada segundo la memoria utilizada, la
utilización del núcleo, la potencia instantánea y la temperatura mediante \texttt{nvidia-smi}. El
tiempo total de ejecución y la memoria residente máxima del proceso se obtienen con
\texttt{/usr/bin/time -v}, y la evolución de la memoria del sistema se registra con
\texttt{vmstat}.

\begin{table}[htbp]
  \centering
  \small
  \caption{Métricas de rendimiento recogidas en cada ejecución y herramienta de medida.}
  \label{tab:metricas}
  \begin{tabular}{lll}
    \toprule
    Métrica & Unidad & Fuente \\
    \midrule
    Precisión del modelo global   & ---            & Evaluación centralizada \\
    Pérdida del modelo global     & ---            & Evaluación centralizada \\
    Tiempo de ejecución           & s              & \texttt{/usr/bin/time -v} \\
    Memoria de GPU (máx. y media) & MiB            & \texttt{nvidia-smi} \\
    Utilización de GPU            & \%             & \texttt{nvidia-smi} \\
    Potencia                      & W              & \texttt{nvidia-smi} \\
    Energía estimada              & Wh             & Cálculo (Ec.~\ref{eq:energia}) \\
    Temperatura de GPU            & \textdegree C  & \texttt{nvidia-smi} \\
    Memoria residente máxima      & KiB            & \texttt{/usr/bin/time -v} \\
    \bottomrule
  \end{tabular}
\end{table}

La energía no se mide con un vatímetro, sino que se estima a partir de la potencia muestreada. Si
se toman $N$ lecturas de potencia $P_i$ a lo largo de una ejecución de duración $T$ segundos, la
energía estimada en vatios-hora se obtiene como el producto de la potencia media por la duración,
expresado en horas:

\begin{equation}
  E_{\text{Wh}} = \frac{1}{3600}\left(\frac{1}{N}\sum_{i=1}^{N} P_i\right) T .
  \label{eq:energia}
\end{equation}

Se trata, por tanto, de una estimación basada en el muestreo de potencia, adecuada para comparar
herramientas entre sí en igualdad de condiciones, pero que no debe interpretarse como una medida
absoluta de consumo eléctrico.

\FloatBarrier
\section{Comparativa de usabilidad: marco y rúbrica}\label{sec:criterios-usabilidad}

La comparativa de usabilidad es cualitativa y persigue valorar qué \textit{framework} resulta más
sencillo de utilizar desde el punto de vista del desarrollador. No se plantea como un estudio de
experiencia de usuario con participantes, sino como una rúbrica de evaluación experta aplicada por
el autor del trabajo, apoyada en evidencias registradas durante la implementación: documentación
consultada, pasos de instalación, errores encontrados, claridad de los registros, facilidad para
configurar los experimentos y capacidad de modificación de cada herramienta.

El marco de valoración se organiza en torno a cuatro dimensiones de la experiencia de uso
—\textit{useful}, \textit{usable}, \textit{desirable} y \textit{accessible}— inspiradas en el
modelo de experiencia de usuario de \textcite{morville2004ux}. Dicho modelo distingue siete
facetas —\textit{useful}, \textit{usable}, \textit{desirable}, \textit{findable},
\textit{accessible}, \textit{credible} y \textit{valuable}—, concebidas para productos de
consumo y, en particular, para sitios web. De ellas se seleccionan y adaptan las cuatro
pertinentes para evaluar una herramienta de desarrollo, que además pueden anclarse en
estándares reconocidos. Las tres restantes se omiten por solaparse, en este contexto, con
las elegidas: \textit{findable} —la facilidad de localizar la documentación, los ejemplos o
las opciones de configuración— queda subsumida en \textit{accessible}; \textit{credible} —la
confianza en la herramienta por su mantenimiento, su comunidad y su estabilidad— se integra
en \textit{desirable}; y \textit{valuable}, que en el modelo de Morville actúa como faceta de
síntesis, coincide en lo esencial con \textit{useful}. La dimensión
\textit{usable} se apoya además en la norma ISO 9241-11, que define la usabilidad en términos de
eficacia, eficiencia y satisfacción en un contexto de uso~\parencite{iso924111}, y en la
descomposición de la usabilidad en aprendizaje, eficiencia, errores y satisfacción del Nielsen
Norman Group~\parencite{nielsen1994usability}. Para la dimensión \textit{accessible} se adapta la
noción de accesibilidad del W3C al contexto de las herramientas software, entendida como facilidad
de acceso a la documentación, los ejemplos, la instalación y los mensajes
comprensibles~\parencite{w3cwcag}. La Tabla~\ref{tab:rubrica-ux} concreta los subcriterios de cada
dimensión.

\begin{table}[htbp]
  \centering
  \small
  \caption{Rúbrica de usabilidad. Cada subcriterio se valora con una escala de 1 a 5.}
  \label{tab:rubrica-ux}
  \begin{tabular}{>{\bfseries}l p{0.52\textwidth} p{0.20\textwidth}}
    \toprule
    \normalfont Dimensión & \normalfont Subcriterios concretos & \normalfont Cómo se mide \\
    \midrule
    Useful &
    Permite implementar FedAvg, usar CIFAR-10 y ResNet-18, registrar métricas, ejecutar varios
    clientes y reproducir experimentos &
    Escala 1--5 según cobertura funcional \\
    \addlinespace
    Usable &
    Instalación, configuración, curva de aprendizaje, claridad de la estructura, facilidad para
    lanzar experimentos y claridad de errores y registros &
    Escala 1--5 según esfuerzo y problemas encontrados \\
    \addlinespace
    Desirable &
    Sensación práctica como desarrollador: comodidad, limpieza de la API, orden del proyecto,
    confianza al trabajar y experiencia de depuración &
    Escala 1--5 con justificación cualitativa \\
    \addlinespace
    Accessible &
    Documentación accesible, ejemplos disponibles, guías actualizadas, mensajes comprensibles,
    requisitos claros y facilidad para empezar sin conocimiento previo profundo &
    Escala 1--5 según barreras de entrada \\
    \bottomrule
  \end{tabular}
\end{table}

El significado de la escala se detalla en la Tabla~\ref{tab:escala-ux}. Para mitigar la
subjetividad inherente a este tipo de valoración, los mismos subcriterios y la misma escala se
aplican por igual a los tres \textit{frameworks}, y cada puntuación se acompaña de la evidencia que
la justifica. Aun así, se reconoce que la valoración refleja el juicio del autor como
desarrollador, lo que se hace constar como limitación en la
Sección~\ref{sec:alcance-limitaciones}.

\begin{table}[htbp]
  \centering
  \small
  \caption{Escala de valoración empleada en la rúbrica de usabilidad.}
  \label{tab:escala-ux}
  \begin{tabular}{cl}
    \toprule
    Puntuación & Significado \\
    \midrule
    1 & Muy deficiente: difícil de usar o con barreras importantes \\
    2 & Deficiente: usable solo con bastante esfuerzo \\
    3 & Aceptable: permite trabajar, pero con limitaciones claras \\
    4 & Bueno: relativamente claro y práctico \\
    5 & Muy bueno: sencillo, claro, completo y bien documentado \\
    \bottomrule
  \end{tabular}
\end{table}

\FloatBarrier
\section{Análisis de flexibilidad: criterios}\label{sec:criterios-flexibilidad}

La flexibilidad se trata como un eje separado de la usabilidad. Mientras que la usabilidad recoge
la percepción de uso, la flexibilidad describe una propiedad arquitectónica: hasta qué punto la
herramienta permite adaptarse a escenarios o necesidades distintos de los previstos por defecto.
Interesa distinguir entre un \textit{framework} cerrado, limitado a las funcionalidades que ya
proporciona, y uno que permite integrar soluciones externas o modificar su funcionamiento.

La valoración se realiza mediante una rúbrica de tres niveles (Tabla~\ref{tab:escala-flex}) sobre
el conjunto de criterios de la Tabla~\ref{tab:criterios-flex}. Estos criterios definen qué se
evalúa en cada caso; la valoración concreta de cada \textit{framework} se presenta en el capítulo
de resultados de usabilidad y flexibilidad.

\begin{table}[htbp]
  \centering
  \small
  \caption{Escala de valoración empleada en el análisis de flexibilidad.}
  \label{tab:escala-flex}
  \begin{tabular}{cl}
    \toprule
    Puntuación & Significado \\
    \midrule
    0 & No soportado o no identificado en la documentación ni en los experimentos \\
    1 & Posible, pero con limitaciones, código manual o soporte experimental \\
    2 & Soporte claro, documentado y relativamente directo \\
    \bottomrule
  \end{tabular}
\end{table}

\begin{table}[htbp]
  \centering
  \small
  \caption{Criterios de flexibilidad evaluados. Cada criterio se valora de 0 a 2.}
  \label{tab:criterios-flex}
  \begin{tabular}{>{\ttfamily}l p{0.30\textwidth} p{0.46\textwidth}}
    \toprule
    \normalfont Código & \normalfont Criterio & \normalfont Qué se evalúa \\
    \midrule
    F1 & Personalización del modelo y el dataset &
    Posibilidad de usar modelos y conjuntos de datos propios, y librerías como PyTorch o
    torchvision, sin rehacer toda la arquitectura. \\
    \addlinespace
    F2 & Personalización del entrenamiento local &
    Facilidad para modificar el bucle de entrenamiento, el optimizador, la función de pérdida, las
    métricas, las épocas locales, el tamaño de lote y la lógica del cliente. \\
    \addlinespace
    F3 & Estrategias de agregación personalizadas &
    Capacidad para implementar o sustituir FedAvg por agregadores propios. \\
    \addlinespace
    F4 & Algoritmos federados disponibles &
    Disponibilidad o facilidad de uso de FedAvg, FedProx, FedOpt, SCAFFOLD u otros algoritmos
    avanzados. \\
    \addlinespace
    F5 & Soporte de aprendizaje federado asíncrono &
    Posibilidad de entrenamiento asíncrono de forma nativa, experimental o solo mediante
    modificaciones profundas. La asincronía reduce las esperas por clientes lentos, pero introduce
    el problema de las actualizaciones obsoletas. \\
    \addlinespace
    F6 & Modos de ejecución y despliegue &
    Soporte de simulación local, ejecución multiproceso, ejecución distribuida y escenarios
    \textit{cross-silo}, \textit{cross-device} o despliegue real. \\
    \addlinespace
    F7 & Integración con herramientas externas &
    Integración con PyTorch, TensorFlow, scikit-learn, Docker, herramientas de registro y
    seguimiento de métricas (MLflow, W\&B) o scripts propios. \\
    \addlinespace
    F8 & Seguridad, privacidad y extensiones avanzadas &
    Posibilidad de añadir privacidad diferencial, agregación segura, filtros, cifrado o políticas
    de seguridad. \\
    \addlinespace
    F9 & Modificación de la arquitectura federada &
    Capacidad de alterar el flujo general: además del esquema servidor-cliente clásico, esquemas
    cíclicos, \textit{swarm}, \textit{split learning} o validación cruzada entre sitios. \\
    \addlinespace
    F10 & Reproducibilidad y configuración &
    Posibilidad de expresar la configuración mediante ficheros, scripts o recetas, y de repetir
    experimentos cambiando pocos parámetros. \\
    \bottomrule
  \end{tabular}
\end{table}

\section{Reproducibilidad y control de variables}\label{sec:reproducibilidad}

Además del entorno descrito en la Sección~\ref{sec:entorno-trabajo}, el estudio adopta varias
medidas orientadas a la reproducibilidad y al control de las variables experimentales.

El control de variables descansa en mantener constantes, entre dos ejecuciones cualesquiera, el
conjunto de datos, el modelo y los hiperparámetros —y un particionamiento IID equivalente entre
herramientas, en el sentido precisado en la Sección~\ref{sec:datos}—, definidos en el
Capítulo~\ref{chap:experimentos}, de manera que las únicas variables sean el \textit{framework}
evaluado y el modo de ejecución. La reproducibilidad se refuerza fijando una semilla aleatoria
común (con valor 42) en los generadores de \texttt{random}, NumPy y PyTorch, y estableciendo
además la semilla de \texttt{hash} del intérprete. Cada configuración se ejecutó una única vez con
esa semilla; por tanto, las métricas se reportan como valores de una sola ejecución y no como
media y desviación típica de varias repeticiones. Esta decisión, motivada por el elevado coste
temporal de los experimentos más largos, se asume como una limitación en la
Sección~\ref{sec:alcance-limitaciones}. Por último, antes de cada ejecución se guarda una
instantánea del entorno —listado de paquetes, exportación del entorno Conda, salida de
\texttt{nvidia-smi} y de \texttt{lscpu}, y revisión del repositorio— que permite trazar las
condiciones exactas de cada corrida. Además, el código de la implementación, los \textit{scripts} de ejecución y los ficheros de configuración de los tres \textit{frameworks} se han publicado en un repositorio público abierto~\parencite{lujan2026repo}, de modo que el estudio puede reproducirse de forma independiente.

\section{Alcance y limitaciones del estudio}\label{sec:alcance-limitaciones}

El estudio compara tres \textit{frameworks} (Flower, NVIDIA FLARE y FedML) sobre un escenario
controlado: CIFAR-10 con partición IID, modelo ResNet-18, algoritmo FedAvg, entre cuatro y diez
clientes simulados en una única máquina, entre diez y treinta rondas, y los modos de ejecución
secuencial y concurrente. La evaluación combina una comparativa cuantitativa de rendimiento con
una comparativa cualitativa de usabilidad y flexibilidad. Quedan fuera del alcance los datos
reales de entornos federados, los clientes distribuidos geográficamente, los algoritmos federados
distintos de FedAvg en la parte cuantitativa y los mecanismos de ataque o defensa de la
privacidad.

La principal limitación es de hardware. La GPU de 6\,GB y los 14\,GiB de memoria RAM condicionaron
el modo concurrente en las tres herramientas, aunque de formas distintas. En Flower, la simulación
se apoya en Ray y el factor limitante fue la memoria RAM del sistema: al alcanzar ciertos umbrales,
Ray terminaba procesos o actores, un comportamiento que se mantuvo con independencia de ajustes
como reducir el número de procesos de carga de datos o liberar la caché. En NVIDIA FLARE, aumentar
la concurrencia mediante el número de hilos implica más procesos y clientes activos a la vez, lo
que resultó limitado por la memoria y el número de procesos en la simulación concurrente; no se
afirma que se produjera un agotamiento de memoria de GPU, al no disponer de un registro que lo
demuestre. En FedML, el backend MPI sí proporcionó concurrencia real entre procesos de cliente,
pero se mostró más sensible al consumo de memoria y podía fallar por falta de recursos en las
configuraciones más grandes. Ningún \textit{framework} quedó, por tanto, libre de impacto en el
modo concurrente, si bien FedML con MPI fue el único que representó una concurrencia real entre procesos.

A esta limitación se añaden otras que conviene tener presentes al interpretar los resultados.
Primero, cada herramienta se ejecutó en un entorno software distinto, con versiones diferentes de
Python, PyTorch y CUDA, lo que constituye una amenaza menor a la validez de la comparación de
rendimiento, mitigada al documentar las versiones exactas. Segundo, los experimentos se limitan a
una distribución IID, por lo que no se evalúa el efecto de la heterogeneidad no~IID. Tercero, los
hiperparámetros se fijaron a priori, sin optimización individual por \textit{framework}, de modo
que los resultados corresponden a esa configuración concreta. Cuarto, cada configuración se
ejecutó una sola vez con una semilla fija, de modo que no se dispone de medidas de variabilidad
entre repeticiones y los valores deben interpretarse como una única observación por configuración.
Quinto, toda la experimentación se
realiza en simulación sobre una sola máquina, por lo que no se miden latencias de red reales ni
efectos propios de un despliegue distribuido. Por último, las métricas de consumo de recursos
dependen del hardware empleado y solo deben extrapolarse con cautela a otros equipos; las
valoraciones de usabilidad y flexibilidad son más generales, pero incorporan el juicio del autor.
